% from ega3-1-fr.tex
\oldpage[III]{112}
prouve donc que l'homomorphisme
\[
\Gamma(X,\mathcal J\otimes\mathcal L^{\otimes n})
\longrightarrow
\Gamma(X,\mathcal J''\otimes\mathcal L^{\otimes n})
\]
est surjectif. Une section $g$ de
$\mathcal J''\otimes\mathcal L^{\otimes n}$ au-dessus de $X$ telle que
$g(x)\neq0$ est donc l'image d'une section
$h\in\Gamma(X,\mathcal J\otimes\mathcal L^{\otimes n})
\subset\Gamma(X,\mathcal L^{\otimes n})$ (car en vertu de
(0\textsubscript{I}, 5.4.1),
$\mathcal J\otimes\mathcal L^{\otimes n}$ est un sous-$\mathcal O_X$-Module
de $\mathcal L^{\otimes n}$)~; on a par définition $h(x)\neq0$ et
$h(z)=0$ pour $z\notin U$, ce qui achève la démonstration.
\end{proof}

\begin{proposition}[2.6.2]
\label{III.2.6.2-fr}
Soient $Y$ un préschéma noethérien, $f:X\to Y$ un morphisme de type
fini, $g:X'\to X$ un morphisme fini surjectif, $\mathcal L$ un
$\mathcal O_X$-Module inversible et $\mathcal L'=g^*(\mathcal L)$. On
suppose vérifiée la condition suivante~: Il existe une partie $Z$ de
$X$, propre sur $Y$ (II, 5.4.10) telle que pour tout $x\in X-Z$, ou bien
$X$ est normal au point $x$, ou bien $(g_*(\mathcal O_{X'}))_x$ est un
$\mathcal O_x$-module libre. Dans ces conditions, pour que $\mathcal L$
soit ample pour $f$, il faut et il suffit que $\mathcal L'$ soit ample
pour $f\circ g$.
\end{proposition}

\begin{proof}
\begin{env}[2.6.2.1]
\label{III.2.6.2.1-fr}
Puisque $g$ est affine, la condition est nécessaire (II, 5.1.12). Pour
voir qu'elle est suffisante, on peut supposer $Y$ affine (II, 4.6.4).
Montrons en outre qu'on peut se borner au cas où $X$ est réduit. En
effet, soit $j:X_{\mathrm{red}}\to X$ l'injection canonique, et posons
$X_1=X_{\mathrm{red}}$, $X'_1=X'\times_X X_1$, de sorte que l'on a le
diagramme commutatif
\[
\xymatrix{
X'\ar[d]_g & X'_1\ar[l]_{j'}\ar[d]^{g_1\, .}\\
X & X_1\ar[l]^{j}
}
\tag{2.6.2.2}\label{III.2.6.2.2-fr}
\]
Le morphisme $f\circ j$ est alors de type fini (I, 6.3.4) et $g_1$ est
un morphisme fini (II, 6.1.5, (iii))~; si $\mathcal L'$ est ample pour
$f\circ g$, $j'^*(\mathcal L')$ est ample pour $f\circ g\circ j'$ puisque
$j'$ est une immersion fermée (II, 5.1.12 et I, 4.3.2). Si on pose
$Z_1=j^{-1}(Z)$, $Z_1$ est propre sur $Y$ (II, 5.4.10)~; d'autre part,
si $X$ est normal en un point $x$, il en est évidemment de même de
$X_{\mathrm{red}}$~; enfin, si $(g_*(\mathcal O_{X'}))_x$ est un
$\mathcal O_x$-module libre, il résulte aussitôt de (II, 1.5.2) que
$((g_1)_*(\mathcal O_{X'_1}))_x$ est un $\mathcal O_{X_1,x}$-module
libre. Enfin, comme $X$ est noethérien (I, 6.3.7), si
$j^*(\mathcal L)$ est ample, $\mathcal L$ est ample (II, 4.5.14), et
comme $j'^*(\mathcal L')=g_1^*(j^*(\mathcal L))$, cela achève la
réduction annoncée. Nous supposons donc désormais $Y$ affine et $X$
réduit.
\end{env}

Les hypothèses de (II, 6.6.11) étant alors vérifiées, il existe un
$Y$-préschéma réduit $X_2$ et un $Y$-morphisme $h:X_2\to X$ fini et
birationnel tel que la restriction de $h$ à $h^{-1}(X-Z)$ soit un
isomorphisme sur $X-Z$ et que $h^*(\mathcal L)$ soit ample. Remplaçant
$X'$ par $X_2$, on voit qu'on est ramené à démontrer la proposition en
supposant en outre que $g$ possède les propriétés qui viennent d'être
énumérées pour $h$. Nous désignerons encore par $Z$ un sous-préschéma de
$X$ ayant $Z$ pour espace sous-jacent, qui est propre sur $Y$
(II, 5.4.10).

\begin{env}[2.6.2.3]
\label{III.2.6.2.3-fr}
Soient maintenant $X_1$ un sous-préschéma fermé de $X$,
$j:X_1\to X$ l'injection canonique,
$X'_1=g^{-1}(X_1)=X'\times_X X_1$ son image réciproque,
$j':X'_1\to X'$ l'injection canonique, de sorte que l'on a le diagramme
commutatif (\hyperref[III.2.6.2.2-fr]{2.6.2.2})~; posons
$\mathcal L_1=j^*(\mathcal L)$,
$\mathcal L'_1=j'^*(\mathcal L')=g_1^*(\mathcal L_1)$, de sorte que
$\mathcal L'_1$ est ample pour $f\circ g\circ j'$ (II, 5.1.12). Si on
pose $Z_1=j^{-1}(Z)$, le sous-préschéma fermé $Z_1$ de $X_1$ est propre
sur $Y$ (II, 5.4.2, (ii)). En d'autres termes, les hypothèses de
(\hyperref[III.2.6.2-fr]{2.6.2}) sont vérifiées pour $X_1$,
$\mathcal L_1$, $g_1$ et $Z_1$.
\end{env}

